\section{Introduction}

Renal cell carcinoma accounts for approximately 90\% of kidney cancers and represents a significant global health burden. Contrast-enhanced computed tomography (CT) is the primary imaging modality for renal tumor assessment, providing essential anatomical detail for diagnosis, surgical planning, and treatment monitoring \cite{ct_imaging_kidney}. Accurate delineation of kidney boundaries and tumor extent is critical for determining resectability, estimating tumor volume, and assessing treatment response.

Manual slice-by-slice segmentation remains the clinical standard but is labor-intensive, time-consuming, and costly. Studies demonstrate substantial inter-observer variability in manual delineation, with different radiologists drawing different boundaries for the same tumor \cite{interobserver_variability}. This variability directly affects surgical planning decisions, treatment monitoring accuracy, and longitudinal follow-up consistency. The subjective nature of manual segmentation introduces uncertainty into clinical workflows and limits reproducibility across institutions.

These limitations motivate automated decision-support tools that provide consistent, reproducible segmentation while preserving physician oversight. Deep learning approaches, particularly three-dimensional convolutional neural networks, have shown promise in medical image segmentation \cite{unet3d}. However, translating these methods into clinically useful systems requires attention to preprocessing, inference strategies, post-processing, and integration with clinical workflows.

This work presents NephroVision, an academic decision-support prototype for automated kidney and renal tumor/cyst segmentation from abdominal CT volumes. The system is trained on the KiTS23 (Kidney Tumor Segmentation Challenge 2023) dataset \cite{kits23}, which provides expert-annotated CT volumes with kidney and tumor/cyst labels. The segmentation pipeline employs a 3D U-Net architecture with Hounsfield-unit preprocessing, sliding-window inference, test-time augmentation, and conservative post-processing. Segmentation masks, volumetric statistics, and interactive three-dimensional visualization are produced for physician review.

Following mid-year feedback emphasizing transparent development processes, the project maintains comprehensive documentation covering preprocessing decisions, training and inference configurations, engineering choices, challenges encountered, and failure analysis. This documentation ensures reproducibility and provides a clear audit trail of design decisions.

The primary contributions are: (1) a complete 3D deep learning pipeline for kidney and renal tumor/cyst segmentation achieving robust kidney segmentation (Dice $0.9307 \pm 0.064$) and reliable tumor detection ($100\%$ on 64 held-out test cases); (2) transparent reporting of moderate tumor boundary precision (Dice $0.6558 \pm 0.262$), reflecting the inherent difficulty of segmenting small, low-contrast lesions; (3) comprehensive development documentation addressing preprocessing, training, engineering choices, challenges, and failure analysis; (4) an interactive web-based visualization system for physician review; and (5) explicit safety framing as an academic decision-support prototype requiring physician oversight, consistent with IEC 62304 Software Safety Class B guidelines \cite{iec62304}.

NephroVision is an academic prototype and is not clinically approved. The system does not claim regulatory clearance and is not intended for autonomous diagnosis or clinical decision-making without physician review. All segmentation outputs require physician verification before clinical use. The system is evaluated exclusively on the KiTS23 dataset; performance on external or multi-center data has not been validated.
